% GENERATED FILE. Edit canonical lessons, then run npm run generate:books.
% canonical-source-hash: fnv1a64:51961ac18dd0022d
% canonical-lessons: ZH-C51-gaoxing, ZH-C51-pa, ZH-C51-leng, ZH-C51-re, ZH-C51-ku

\chapter{How You Feel}
\label{ch:zh-how-you-feel}
\hlchaptermodality{51}

\begin{chapteropening}
\textbf{By the end of this chapter:} \emph{I can say I am happy, afraid, cold or hot.}

Complete the last lesson of chapter 51: I can say I am happy, afraid, cold or hot.
\end{chapteropening}

\section[gāoxìng]{gāoxìng — happy; glad}

\label{lesson:ZH-C51-gaoxing}



\begin{quote}
Before the new one: say the Mandarin for health, then the Mandarin for to fall ill.
\end{quote}

\subsection*{You'll want to know: gāoxìng}
\emph{gāoxìng} — \textquotedblleft{}happy; glad\textquotedblright{}.

\begin{quote}
\emph{wǒ hěn gāoxìng} — I'm so pleased.
\end{quote}

Its characters come later; for now it is a word for the ear.

\begin{cousinweb}
\emph{gāoxìng} is built from gāo 'high' and xìng 'mood', both inherited from Middle Chinese, \emph{kaw} and \emph{xingH} (Guangyun, in Baxter's transcription). The dictionary glosses the first 'high'; it gives the second reading no gloss of its own, so 'mood' is the usual explanation, not a settled one.
\end{cousinweb}

\begin{grammarlens}[title={What you've built}]
One more word for saying how you are.
\end{grammarlens}

\subsection*{Your turn}
\begin{itemize}
\raggedright
  \item \emph{Say it:} \emph{gāoxìng}
  \item \emph{Say it:} \emph{gāoxìng}, once more
  \item \emph{Say it:} \emph{wǒ hěn gāoxìng} — I'm so pleased.
  \item \emph{Recall:} say the Mandarin for health, then the Mandarin for to fall ill, then say \emph{gāoxìng} again
\end{itemize}

\begin{tcolorbox}[breakable,colback=teal!4,colframe=teal!35!black,title={Before you move on}]
What does gāoxìng mean? (\textquotedblleft{}Happy; glad\textquotedblright{}.) Say it once more.
\end{tcolorbox}
\section[pà]{pà — to be afraid of; to fear}

\label{lesson:ZH-C51-pa}



\begin{quote}
Before the new one: say the Mandarin for to fall ill, then the Mandarin for happy.
\end{quote}

\subsection*{You'll want to know: pà}
\emph{pà} — \textquotedblleft{}to be afraid of; to fear\textquotedblright{}.

\begin{quote}
\emph{nǐ pà shénme} — What are you afraid of?
\end{quote}

Its characters come later; for now it is a word for the ear.

\begin{cousinweb}
\emph{pà} is inherited from Middle Chinese \emph{phaeH} (Guangyun, in Baxter's transcription), glossed 'to fear'.
\end{cousinweb}

\begin{grammarlens}[title={What you've built}]
One more word for saying how you are.
\end{grammarlens}

\subsection*{Your turn}
\begin{itemize}
\raggedright
  \item \emph{Say it:} \emph{pà}
  \item \emph{Say it:} \emph{pà}, once more
  \item \emph{Say it:} \emph{nǐ pà shénme} — What are you afraid of?
  \item \emph{Recall:} say the Mandarin for to fall ill, then the Mandarin for happy, then say \emph{pà} again
\end{itemize}

\begin{tcolorbox}[breakable,colback=teal!4,colframe=teal!35!black,title={Before you move on}]
What does pà mean? (\textquotedblleft{}To be afraid of; to fear\textquotedblright{}.) Say it once more.
\end{tcolorbox}
\section[lěng]{lěng — cold}

\label{lesson:ZH-C51-leng}



\begin{quote}
Before the new one: say the Mandarin for happy, then the Mandarin for to be afraid.
\end{quote}

\subsection*{You'll want to know: lěng}
\emph{lěng} — \textquotedblleft{}cold\textquotedblright{}.

\begin{quote}
\emph{jīntiān hěn lěng} — It's cold today.
\end{quote}

Its characters come later; for now it is a word for the ear.

\begin{cousinweb}
\emph{lěng} is inherited from Middle Chinese \emph{lengX} (Guangyun, in Baxter's transcription), glossed 'cold'.
\end{cousinweb}

\begin{grammarlens}[title={What you've built}]
One more word for saying how you are.
\end{grammarlens}

\subsection*{Your turn}
\begin{itemize}
\raggedright
  \item \emph{Say it:} \emph{lěng}
  \item \emph{Say it:} \emph{lěng}, once more
  \item \emph{Say it:} \emph{jīntiān hěn lěng} — It's cold today.
  \item \emph{Recall:} say the Mandarin for happy, then the Mandarin for to be afraid, then say \emph{lěng} again
\end{itemize}

\begin{tcolorbox}[breakable,colback=teal!4,colframe=teal!35!black,title={Before you move on}]
What does lěng mean? (\textquotedblleft{}Cold\textquotedblright{}.) Say it once more.
\end{tcolorbox}
\section[rè]{rè — hot}

\label{lesson:ZH-C51-re}



\begin{quote}
Before the new one: say the Mandarin for to be afraid, then the Mandarin for cold.
\end{quote}

\subsection*{You'll want to know: rè}
\emph{rè} — \textquotedblleft{}hot\textquotedblright{}.

\begin{quote}
\emph{zhōngwǔ hěn rè} — It's hot at midday.
\end{quote}

Its characters come later; for now it is a word for the ear.

\begin{cousinweb}
\emph{rè} is inherited from Middle Chinese \emph{nyet} (Guangyun, in Baxter's transcription), which the dictionary explains as burning, as if by fire.
\end{cousinweb}

\begin{grammarlens}[title={What you've built}]
One more word for saying how you are.
\end{grammarlens}

\subsection*{Your turn}
\begin{itemize}
\raggedright
  \item \emph{Say it:} \emph{rè}
  \item \emph{Say it:} \emph{rè}, once more
  \item \emph{Say it:} \emph{zhōngwǔ hěn rè} — It's hot at midday.
  \item \emph{Recall:} say the Mandarin for to be afraid, then the Mandarin for cold, then say \emph{rè} again
\end{itemize}

\begin{tcolorbox}[breakable,colback=teal!4,colframe=teal!35!black,title={Before you move on}]
What does rè mean? (\textquotedblleft{}Hot\textquotedblright{}.) Say it once more.
\end{tcolorbox}
\section[kū]{kū — to cry; to weep}

\label{lesson:ZH-C51-ku}



\begin{quote}
Before the new one: say the Mandarin for cold, then the Mandarin for hot.
\end{quote}

\subsection*{You'll want to know: kū}
\emph{kū} — \textquotedblleft{}to cry; to weep\textquotedblright{}.

\begin{quote}
\emph{háizi kū le} — The child started crying.
\end{quote}

Its characters come later; for now it is a word for the ear.

\begin{cousinweb}
\emph{kū} is inherited from Middle Chinese \emph{khuwk} (Guangyun, in Baxter's transcription), glossed 'a sound of grief'.
\end{cousinweb}

\begin{grammarlens}[title={What you've built}]
One more word for saying how you are.
\end{grammarlens}

\subsection*{Your turn}
\begin{itemize}
\raggedright
  \item \emph{Say it:} \emph{kū}
  \item \emph{Say it:} \emph{kū}, once more
  \item \emph{Say it:} \emph{háizi kū le} — The child started crying.
  \item \emph{Recall:} say the Mandarin for cold, then the Mandarin for hot, then say \emph{kū} again
\end{itemize}

\begin{tcolorbox}[breakable,colback=teal!4,colframe=teal!35!black,title={Before you move on}]
What does kū mean? (\textquotedblleft{}To cry; to weep\textquotedblright{}.) Say it once more.
\end{tcolorbox}
